% 对应题号 M2-045，本题条件、变量及几何关系逐项核对。
% 修订：启用中文断行；核对题设及图文对应。
\documentclass[tikz,border=14pt]{standalone}
\usepackage{fontspec,xeCJK}
\usepackage{amsmath,amssymb}
\usepackage{xcolor}
\usepackage{tikz}
\usetikzlibrary{arrows.meta,calc,patterns,positioning,bending}
\IfFontExistsTF{Heiti SC}{\IfFontExistsTF{Heiti SC}{\setmainfont{Heiti SC}}{\setmainfont{FandolHei-Regular}}\IfFontExistsTF{Heiti SC}{\setsansfont{Heiti SC}}{\setsansfont{TeX Gyre Heros}}}{\IfFontExistsTF{Songti SC}{\setmainfont{Songti SC}}{\setmainfont{FandolSong-Regular}}\IfFontExistsTF{Songti SC}{\setsansfont{Songti SC}}{\setsansfont{FandolSong-Regular}}}
\definecolor{bluefield}{RGB}{24,76,255}
\definecolor{redfield}{RGB}{232,49,35}
\definecolor{orangefield}{RGB}{236,132,18}
\definecolor{greenfield}{RGB}{18,132,86}
\definecolor{purplefield}{RGB}{112,72,190}
\definecolor{linegray}{RGB}{60,68,78}
\definecolor{softblue}{RGB}{238,243,255}
\definecolor{softgreen}{RGB}{238,249,244}
\definecolor{softorange}{RGB}{255,247,235}
\tikzset{
  arrow/.style={-{Latex[length=3.1mm,width=2.1mm]},line width=1.1pt},
  thinarr/.style={-{Latex[length=2.5mm,width=1.7mm]},line width=.85pt},
  axis/.style={-{Latex[length=2.5mm,width=1.7mm]},line width=.8pt,draw=linegray},
  guide/.style={draw=linegray!55,dashed,line width=.65pt},
  box/.style={draw=linegray,rounded corners=3pt,fill=white,align=center,inner sep=6pt,line width=.75pt},
  bluebox/.style={box,draw=bluefield,fill=softblue},
  greenbox/.style={box,draw=greenfield,fill=softgreen},
  orangebox/.style={box,draw=orangefield,fill=softorange}
}
\IfFontExistsTF{Heiti SC}{\setCJKmainfont[AutoFakeBold=1.5]{Heiti SC}}{\setCJKmainfont[AutoFakeBold=1.5]{FandolHei-Regular}}
\IfFontExistsTF{Heiti SC}{\setCJKsansfont{Heiti SC}}{\setCJKsansfont{FandolHei-Regular}}
\begin{document}

\begin{tikzpicture}[x=1cm,y=1cm]
\path[use as bounding box] (-10,-4.6) rectangle (10,4.6);
\fill[white] (-10,-4.6) rectangle (10,4.6);
\node[font=\Large\bfseries,text width=18.6cm,align=center] at (0,4.05) {轴对称偶极与四极场线：角向节点决定特殊场线};
\begin{scope}[shift={(-5.0,-.15)}]
\node[font=\bfseries] at (0,2.75) {$L=1$}; \fill[linegray] (0,0) circle(.06);
\foreach \s in {1.1,1.7,2.3}{\draw[bluefield] plot[domain=8:172,samples=70] ({\s*sin(\x)^3},{\s*sin(\x)^2*cos(\x)}); \draw[bluefield] plot[domain=188:352,samples=70] ({\s*sin(\x)^3},{\s*sin(\x)^2*cos(\x)});}
\node[bluefield] at (0,-2.6) {$r=C\sin^2\theta$};
\end{scope}
\begin{scope}[shift={(0,-.15)}]
\node[font=\bfseries] at (0,2.75) {$L=2$}; \draw[axis] (0,-2.2)--(0,2.2); \draw[axis] (-2.2,0)--(2.2,0);
\foreach \s in {.8,1.2,1.6}{\foreach \sx in {-1,1}{\foreach \sy in {-1,1}{\draw[greenfield] plot[domain=10:80,samples=45] ({\sx*\s*sqrt(sin(\x)^2*cos(\x))*sin(\x)},{\sy*\s*sqrt(sin(\x)^2*cos(\x))*cos(\x)});}}}
\node[greenfield] at (0,-2.6) {$r^2=C\sin^2\theta\cos\theta$};
\end{scope}
\begin{scope}[shift={(5.2,-.1)}]
\node[bluebox,text width=5.6cm] at (0,1.45) {$\displaystyle \frac{dr}{r\,d\theta}=\frac{(L+1)P_L}{\sin\theta\,P_L'}$};
\node[greenbox,text width=5.6cm] at (0,-.1) {积分时单独检查轴线、赤道面及 $P_L'=0$ 的方向。};
\node[orangebox,text width=5.2cm,align=center] at (0,-1.75) {理想多极公式只描述远区；\\不能替代有限源附近的精确场线。};
\end{scope}
\end{tikzpicture}
\end{document}
